\documentclass[11pt,letterpaper]{article}
\usepackage[margin=1in]{geometry}
\usepackage{graphicx}
\usepackage{hyperref}
\usepackage{xcolor}
\usepackage{listings}
\usepackage{amsmath}
\usepackage{booktabs}
\usepackage{float}
\usepackage{fancyhdr}
\usepackage{caption}

\pagestyle{fancy}
\fancyhf{}
\rhead{EEL4783 Final Project Part 1}
\lhead{Sobel Edge Detection --- HLS}
\rfoot{Page \thepage}

\lstset{
  basicstyle=\ttfamily\footnotesize,
  breaklines=true,
  frame=single,
  language=Verilog,
  numbers=left,
  numberstyle=\tiny\color{gray},
  tabsize=2
}

\title{\textbf{EEL4783 Final Project (Part 1/2):\\High-Level Synthesis of a Sobel Edge Detection Filter}}
\author{
  % TODO: Replace with your name and NID
  Your Name \\
  NID: yo485591 \\
  University of Central Florida \\
  EEL4783 --- HDL in Digital Systems
}
\date{April 15, 2026}

\begin{document}
\maketitle
\thispagestyle{fancy}

% ============================================================
\section*{Introduction}
% ============================================================

This report documents the High-Level Synthesis (HLS) of a Sobel Edge Detection filter using Vivado HLS 2019.1.
The design takes a 256$\times$256 grayscale image as input and produces an edge-detected image using the Sobel operator.
The Sobel filter applies 3$\times$3 horizontal ($G_x$) and vertical ($G_y$) convolution kernels and computes the gradient magnitude
$|\nabla I| \approx |G_x| + |G_y|$, clamped to the pixel range $[0, 255]$.

The design was synthesized with default settings (no HLS pragmas) targeting a Zynq-7020 (\texttt{xc7z020clg484-1}) with a 10\,ns clock period.
The following sections present the synthesis results, C/RTL co-simulation, and generated RTL.

% ============================================================
\section{Question 1: Synthesis Results}
% ============================================================

\subsection{1a) Default Synthesis Metrics}

The design was synthesized using the default Vivado HLS scheduling (no optimization pragmas).
Table~\ref{tab:synth} summarizes the key metrics from the synthesis report.

\begin{table}[H]
\centering
\caption{Default Synthesis Results (Sobel Edge Detection Filter)}
\label{tab:synth}
\begin{tabular}{@{}ll@{}}
\toprule
\textbf{Metric}             & \textbf{Value} \\
\midrule
Estimated Clock Period       & \textit{TODO: e.g., 8.70\,ns} \\
Worst-Case Latency (cycles)  & \textit{TODO: e.g., 2,325,143} \\
Worst-Case Latency (time)    & \textit{TODO: e.g., 23.25\,ms} \\
Number of DSP48E Used        & \textit{TODO: e.g., 2} \\
Number of FFs Used           & \textit{TODO: e.g., 200} \\
Number of LUTs Used          & \textit{TODO: e.g., 500} \\
\bottomrule
\end{tabular}
\end{table}

% TODO: Insert screenshot of the synthesis summary from Vivado HLS
% \begin{figure}[H]
%   \centering
%   \includegraphics[width=0.9\textwidth]{screenshots/synth_summary.png}
%   \caption{Vivado HLS Synthesis Summary for sobel\_filter.}
%   \label{fig:synth_summary}
% \end{figure}

The estimated clock period is below the 10\,ns target constraint, confirming the design meets timing.
The design is not pipelined (0 out of 4 loops pipelined), which means the worst-case latency equals the total
number of sequential clock cycles needed to process all $256 \times 256 = 65{,}536$ pixels through the
nested convolution loops.

\subsection{1b) Analysis Perspective}

Opening the Analysis Perspective (Solution $\rightarrow$ Open Analysis Perspective) provides a cycle-level
view of how operations are scheduled across the FSM states.

% TODO: Insert screenshot of the Analysis Perspective / Schedule Viewer
% \begin{figure}[H]
%   \centering
%   \includegraphics[width=0.9\textwidth]{screenshots/analysis_perspective.png}
%   \caption{Analysis Perspective showing the schedule of operations within the sobel\_filter FSM states.}
%   \label{fig:analysis}
% \end{figure}

The schedule viewer reveals that each iteration of the innermost kernel loop (the 3$\times$3 convolution)
requires multiple states for memory reads from the \texttt{src} array, multiply-accumulate operations
with the Sobel kernel coefficients, and the final magnitude computation and clamped write to \texttt{dst}.

\subsection{1c) Resource Profile}

The Resource Profile tab shows how hardware resources are distributed across the design.

% TODO: Insert screenshot of the Resource Profile tab
% \begin{figure}[H]
%   \centering
%   \includegraphics[width=0.9\textwidth]{screenshots/resource_profile.png}
%   \caption{Resource Profile showing DSP48E, FF, and LUT allocation within the sobel\_filter module.}
%   \label{fig:resource}
% \end{figure}

The Bind Op Report confirms that the DSP48E blocks are allocated to the multiply-accumulate operations
within the inner kernel loop, where each pixel is multiplied by its corresponding Sobel kernel coefficient.
The Bind Storage Report shows that the Sobel kernels ($G_x$ and $G_y$) are implemented as distributed ROMs
(\texttt{sobel\_filter\_Gx\_rom} and \texttt{sobel\_filter\_Gy\_rom}), consistent with the synthesis log output.

The input and output image arrays (\texttt{src} and \texttt{dst}) are accessed through external \texttt{ap\_memory}
interfaces rather than on-chip BRAM, as expected for arrays of $65{,}536$ bytes each.

% ============================================================
\section{Question 2: C/RTL Co-simulation}
% ============================================================

\subsection{2a) Co-simulation Report}

C-simulation was first run to confirm functional correctness of the C++ source and testbench.
The testbench generates a synthetic 256$\times$256 image with vertical stripe patterns and a horizontal bar,
applies the Sobel filter, and compares the output against a software reference implementation pixel-by-pixel.

% TODO: Insert screenshot of C-simulation output
% \begin{figure}[H]
%   \centering
%   \includegraphics[width=0.8\textwidth]{screenshots/csim_output.png}
%   \caption{C-simulation completes with 0 errors, confirming the C++ source and testbench are functionally correct.}
%   \label{fig:csim}
% \end{figure}

The C-simulation output confirms:
\begin{verbatim}
Non-zero output pixels: 4264 / 65536
TEST PASSED (4264 edge pixels detected)
\end{verbatim}

C/RTL co-simulation was then run using the Verilog RTL model under XSIM.
The RTL simulation feeds the same testbench stimulus into the synthesized Verilog and compares
the outputs against the software reference.

% TODO: Insert screenshot of co-simulation report
% \begin{figure}[H]
%   \centering
%   \includegraphics[width=0.8\textwidth]{screenshots/cosim_report.png}
%   \caption{C/RTL co-simulation report showing PASS with measured latency.}
%   \label{fig:cosim}
% \end{figure}

The co-simulation completed with status \textbf{PASS}, confirming that the synthesized RTL produces
identical results to the C model. The simulation log reports:

\begin{verbatim}
RTL Simulation : 1 / 1 [100.00%] @ "23251435000"
$finish called at time : 23251475 ns
\end{verbatim}

The measured RTL latency of approximately 23.25\,ms (at 10\,ns clock period) is consistent with
the synthesis estimate, confirming that the un-pipelined design processes the full 256$\times$256 image
sequentially through the nested convolution loops.

Since only one transaction was simulated, the initiation interval (II) is reported as N/A.
For this non-pipelined design, the next input can be provided after latency${}+1$ clock cycles.

\subsection{2b) Waveform Trace Analysis}

The co-simulation dumps waveform traces that can be viewed in Vivado's waveform viewer.

% TODO: Insert screenshot of the waveform traces
% \begin{figure}[H]
%   \centering
%   \includegraphics[width=0.9\textwidth]{screenshots/cosim_waveform.png}
%   \caption{Co-simulation waveform showing the ap\_ctrl\_hs handshake signals (AP\_CLK, AP\_START, AP\_READY, AP\_DONE).}
%   \label{fig:waveform}
% \end{figure}

The waveform trace confirms the standard \texttt{ap\_ctrl\_hs} protocol behavior:
\begin{itemize}
  \item \texttt{AP\_START} is asserted high by the testbench to initiate processing.
  \item The design processes the full image over the measured latency period.
  \item \texttt{AP\_DONE} pulses high once all output pixels have been written.
  \item \texttt{AP\_READY} indicates the design can accept a new input.
  \item Memory interface signals (\texttt{src\_address0}, \texttt{src\_ce0}, \texttt{dst\_we0}) show
        the sequential read/write pattern of the un-pipelined processing loop.
\end{itemize}

\subsection{2c) Export RTL}

The synthesis step generates both VHDL and Verilog RTL under the solution directory:
\begin{itemize}
  \item \texttt{sobel\_hls/solution1/syn/verilog/} --- Verilog sources
  \item \texttt{sobel\_hls/solution1/syn/vhdl/} --- VHDL sources
\end{itemize}

\noindent The generated RTL consists of:
\begin{itemize}
  \item \textbf{sobel\_filter.v}: Top-level module containing the FSM, address generation,
        loop counters, and pipeline control logic.
  \item \textbf{sobel\_filter\_Gx.v / sobel\_filter\_Gy.v}: Distributed ROM modules storing
        the $3 \times 3$ Sobel kernel coefficients.
  \item \textbf{sobel\_filter\_mac\_muladd\_8ns\_3s\_32ns\_32\_1\_1.v}: Multiply-accumulate
        primitive mapped to DSP48E blocks, used for the kernel convolution.
\end{itemize}

\textbf{Note:} The IP Catalog export step (\texttt{export\_design}) fails on Vivado HLS 2019.1
due to a known timestamp overflow bug (\texttt{core\_revision} exceeds 32-bit integer range after 2024).
The generated RTL is fully available in the synthesis output directory and is functionally complete.

% ============================================================
\section*{Closing Notes}
% ============================================================

The default synthesis run produces a working Sobel Edge Detection filter that processes
a 256$\times$256 grayscale image on the Zynq-7020 target.
The design uses \texttt{ap\_memory} interfaces for the input and output image arrays,
with the Sobel kernels implemented as distributed ROMs.

Without any optimization pragmas, the design processes the image sequentially through
four nested loops (row, column, kernel-y, kernel-x), resulting in high latency.
Part~2 of this project will explore the design space by applying HLS pragmas
(loop pipelining, array partitioning, loop unrolling) to reduce latency
and improve throughput at the cost of additional resource usage.

% ============================================================
\newpage
\appendix
\section{Appendix A: HLS C++ Source Code}
% ============================================================

\subsection{A.1 sobel.h}
\lstinputlisting[language=C++]{sobel.h}

\subsection{A.2 sobel.cpp}
\lstinputlisting[language=C++]{sobel.cpp}

\subsection{A.3 sobel\_tb.cpp (Testbench)}
\lstinputlisting[language=C++]{sobel_tb.cpp}

% ============================================================
\section{Appendix B: Generated Verilog Source}
% ============================================================

% TODO: After running synthesis on the VLSI server, copy the generated Verilog files
% to a local directory and uncomment the following:
%
% \subsection{B.1 sobel\_filter.v (Top Level)}
% \lstinputlisting[language=Verilog]{verilog/sobel_filter.v}
%
% \subsection{B.2 sobel\_filter\_Gx.v}
% \lstinputlisting[language=Verilog]{verilog/sobel_filter_Gx.v}
%
% \subsection{B.3 sobel\_filter\_Gy.v}
% \lstinputlisting[language=Verilog]{verilog/sobel_filter_Gy.v}
%
% \subsection{B.4 sobel\_filter\_mac\_muladd\_8ns\_3s\_32ns\_32\_1\_1.v}
% \lstinputlisting[language=Verilog]{verilog/sobel_filter_mac_muladd_8ns_3s_32ns_32_1_1.v}

\end{document}
